%% annex/p09/proofs.tex --- editorial proofs for Part IX

\begin{annexitem}{annex:p9.alc.1}{thm:p09-alc-tableau}{Correctness and termination of the \(\mathcal{ALC}\) completion procedure}
Put
\(\mathcal C=\operatorname{cl}(C_0,\mathcal T,\mathcal A)\), as in
Definition~\ref{def:p09-alc-tableau}, and let \(N=|\mathcal C|\).  Every label
is a subset of the finite set \(\mathcal C\).  A Boolean, universal, or global
step strictly enlarges one label; each disjunction occurrence is resolved once
on a branch; and each existential formula at a node creates at most one
witness because its witness persists.  Only the existential rule creates
nodes, so the number of witness children of a node is bounded by the number of
existential formulas in \(\mathcal C\); the initial roots and declared ABox
edges are finite fixed input data.  At every stage, an unblocked
witness-tree root-to-node path has no descendant label contained in an
ancestor label, and hence has at most \(2^N\) unblocked nodes.  A new witness
can be created only at an unblocked source, so no witness-tree node can occur
below depth \(2^N\).  The finite depth and finite witness branching bound the
number of completion nodes.  Each such node permits only finitely many strict
label additions and branch choices.  Every legal expansion sequence therefore
terminates.

For completeness, suppose that an interpretation \(I\) satisfies the TBox and
ABox and, when \(C_0\) is present, has some element in \(C_0^I\).  Map every
named root \(a\) to \(a^I\), the tested root to such an element, and a
TBox-only anonymous root to
any element of the nonempty domain.  The initialization rule satisfies
\[
 E\in\mathcal L(x)\Longrightarrow\pi(x)\in E^I,
 \qquad
 x\mathrel r y\Longrightarrow(\pi(x),\pi(y))\in r^I:
\]
ABox concept assertions give the first clause at named roots, ABox role
assertions give the second, the tested concept gives the first at \(c_0\), and
every global internalization is true at every domain element.  Conjunction and
universal propagation preserve the invariant.  At a disjunction choose a true
disjunct, and at an existential choose a semantic witness as the image of the
fresh node.  Blocking suppresses a construction step but adds no formula, so
it cannot create a clash.  Exhausting these model-guided choices yields a
complete clash-free branch.

For soundness, fix a complete clash-free branch and let the interpretation
domain be the unblocked completion nodes.  Redirection is defined only where
the role interpretation needs it.  For every retained edge
\(x\mathrel r y\) whose source \(x\) is unblocked, put
\[
 \widehat y=
 \begin{cases}
  y,&y\text{ is unblocked},\\
  b(y),&y\text{ is blocked},
 \end{cases}
\]
where in the second case \(b(y)\) is a directly blocking ancestor satisfying
\(\mathcal L(y)\subseteq\mathcal L(b(y))\).  This choice is defined on every
such target: if \(y\) is blocked while its immediate predecessor \(x\) is
unblocked, no ancestor of \(y\) is blocked, so \(y\) is directly blocked and
each possible blocker is unblocked.  Residual nodes blocked only through an
ancestor are outside the domain of this redirection.  Interpret a concept name
\(A\) by
\[
 A^I=\{x:x\text{ is unblocked and }A\in\mathcal L(x)\},
\]
and put \((x,\widehat y)\in r^I\) for every retained edge \(x\mathrel r y\) whose
source \(x\) is unblocked.  Named roots are never blocked, so interpret each
individual by its root.

For every unblocked \(x\) and \(E\in\mathcal L(x)\), structural induction on
\(E\) proves \(x\in E^I\).  Clash-freeness handles concept names and their NNF
negations; completion handles conjunction and disjunction.  If
\(\exists r.D\in\mathcal L(x)\), completeness and the fact that \(x\) is
unblocked supply an edge \(x\mathrel r y\) with
\(D\in\mathcal L(y)\).  If \(y\) is blocked, subset blocking gives
\(\mathcal L(y)\subseteq\mathcal L(\widehat y)\), so the induction hypothesis applies
at the redirected target.  If \(\forall r.D\in\mathcal L(x)\), universal
propagation puts \(D\) in every original target label; the same subset argument
transfers it to each redirected target.  Hence the interpretation satisfies
every labelled concept.

Every global internalization labels every node, so the interpretation satisfies
\(\mathcal T\).  Initialized named-root labels and retained declared edges
satisfy \(\mathcal A\), and the tested root satisfies \(C_0\) when present.
Thus a complete clash-free branch yields a finite model.  The model-guided and
model-extraction directions, together with termination, prove
Theorem~\ref{thm:p09-alc-tableau}.
\end{annexitem}

\begin{annexitem}{annex:p9.ctl.1}{thm:p09-ctl-labelling}{Correctness and complexity of CTL labelling}
Proceed by structural induction on the processed CTL formula.  Atomic,
negation, conjunction, and disjunction clauses are exactly their truth-set
operations.  The two predecessor clauses are the semantic clauses for
\(\mathsf{EX}\) and \(\mathsf{AX}\).

For the least fixed point defining \(\mathsf E[B\mathsf U C]\), iteration from
\(\varnothing\) has the following stage invariant: a state enters by stage
\(i+1\) exactly when it begins a path of at most \(i\) transitions whose final
state satisfies \(C\) and whose earlier states satisfy \(B\).  The claim follows
by induction on \(i\), using \(\operatorname{Pre}_{\exists}\) for the first
edge.  Taking the union of the stages gives precisely the existential-until
semantics.

For \(\mathsf A[B\mathsf U C]\), the analogous stage invariant says that every
path reaches a \(C\)-state within the common bound \(i\), with \(B\) before
that point.  The induction uses \(\operatorname{Pre}_{\forall}\).  If the CTL
semantics holds but no finite common bound exists, finite branching supplies,
by K\"onig's lemma, an infinite path that remains in \(B\setminus C\), a
contradiction.  Thus the least fixed point is exactly universal until.

Iteration for \(\mathsf{EG}B\) from \(S\) removes precisely states from which
no arbitrarily long \(B\)-path begins.  In a finite graph, arbitrarily long
such paths yield an infinite one; conversely the states of any infinite
\(B\)-path form a post-fixed set.  This proves the greatest-fixed-point clause.
The \(\mathsf{AG}B\) clause is the universal analogue: its greatest fixed point
is the largest successor-closed subset of \(L(B)\).  The derived
\(\mathsf{EF}\) and \(\mathsf{AF}\) cases are the until cases with left operand
\(\top\).  This completes the truth-set induction.

For complexity, store one bit per state and subformula.  Boolean and predecessor
clauses scan states and adjacency lists once.  In an until or globally clause,
a worklist charges each state insertion or deletion once and each incident edge
when that change is propagated.  Hence each subformula costs
\(O(|S|+|R|)\), and all subformulas cost
\(O((|S|+|R|)|A|)\), as claimed.
\end{annexitem}

\begin{annexitem}{annex:p9.hoare.1}{thm:p09-hoare-soundness}{Soundness of the registered Hoare rules}
Induct on a derivation on
\hyperref[sheet:hoare-while]{\texttt{sheet:hoare-while}}.  Skip preserves the
state.  Assignment followed from a state satisfying \(Q[e/x]\) produces
exactly the updated state satisfying \(Q\); this is the substitution lemma for
total expressions.  Consequence composes the two valid assertion implications
with the induction hypothesis.  Sequence composes the intermediate state of
the first big step with the second.  A conditional big step takes exactly the
branch selected by \(b\), so the corresponding premise applies.  Each of these
arguments preserves both the partial and total tag; in the total case the
induction hypotheses also supply termination of the premise commands.

For the partial while rule, induct on the number of body executions in a
terminating big-step derivation.  With zero executions, \(I\land\neg b\) holds.
Otherwise \(I\land b\) holds initially, the premise returns a state satisfying
\(I\), and the induction hypothesis applies to the remaining loop derivation.

For the total while rule, first take an arbitrary state satisfying \(I\).  If
\(b\) is false, the loop exits immediately, with the required postcondition.
If \(b\) is true, the first side premise makes the old integer value
\(n=\llbracket t\rrbracket\) nonnegative.  Prove termination from every such
guarded state by well-founded induction on \(n\in\mathbb N\).  Give the fresh
rigid variable \(z\) the value \(n\).  The total-correctness premise terminates
the body in a state satisfying \(I\) and having variant value \(n'<n\).  If
the guard is now false, the loop exits.  If it remains true, the first side
premise gives \(n'\geq0\), so the induction hypothesis applies to the
remaining loop.  Thus every initial \(I\)-state terminates, whether it exits
before the induction is needed or enters with a natural-valued measure.  Its
final state satisfies \(I\land\neg b\).  Every registered rule is therefore
sound for its printed correctness tag.
\end{annexitem}

\begin{annexitem}{annex:p9.tqbf.1}{thm:p04-tqbf-pspace}{TQBF completeness: membership and polynomial-space simulation}
Membership is direct.  Evaluate a closed quantified Boolean formula by a
depth-first recursion on its syntax.  At \(\exists p\), evaluate the body first
with \(p=0\) and, if necessary, with \(p=1\); at \(\forall p\), require both
values.  The two branches are explored sequentially.  A current assignment,
syntax position, and recursion stack use space polynomial in the formula
length, so \(\mathsf{TQBF}\in\mathsf{PSPACE}\).

For hardness, let deterministic machine \(M\) use at most \(p(n)\) work cells.
Modify it, without changing its language or polynomial-space bound, so that
acceptance first erases the work tapes, resets the encoded heads, and enters
one canonical accepting configuration.  Encode each configuration on input
\(x\) by an \(m\)-bit block, where \(m\geq1\) is polynomial in \(|x|\).
The read-only input contents are fixed by \(x\), rather than copied into the
work-tape block.
There is a polynomial-size Boolean formula
\(\operatorname{Step}_{M,x}(u,v)\) which holds exactly when valid block \(v\)
is the unique legal successor of valid block \(u\): it checks the state, head,
and each tape cell by the finite local transition table, and it is false when
either block is invalid.  Write
\(\operatorname{Eq}(u,v)\) for bitwise equality of two \(m\)-bit blocks.

The following are metalinguistic formula abbreviations, introduced solely to
prevent the usual duplicated-recursion blow-up.  Their standing is a family of
quantified Boolean formulas with two free \(m\)-bit parameter blocks; they add
no connective to QBF.  Put
\[
 R_0(u,v)=\operatorname{Eq}(u,v)\lor
           \operatorname{Step}_{M,x}(u,v).
\]
For fresh \(m\)-bit blocks \(z,a,c\) and a fresh bit \(b\), let
\[
\begin{aligned}
 \operatorname{Sel}_b(u,z,v;a,c)
  ={}&\bigl(\neg b\to
       (\operatorname{Eq}(a,u)\land\operatorname{Eq}(c,z))\bigr)\\
    &\land\bigl(b\to
       (\operatorname{Eq}(a,z)\land\operatorname{Eq}(c,v))\bigr),\\
 R_{i+1}(u,v)
  ={}&\exists z\,\forall b\,\exists a\,\exists c\,
       \bigl(\operatorname{Sel}_b(u,z,v;a,c)\land R_i(a,c)\bigr).
\end{aligned}
\]
Here a block quantifier abbreviates \(m\) ordinary Boolean quantifiers.  The
selector is necessary because it makes the universally chosen half use one,
not two, textual occurrences of \(R_i\).  Its quantified blocks \(a,c\) may
depend on \(b\).  Consequently
\[
 |R_{i+1}|=|R_i|+O(m),
\]
apart from the fixed polynomial-size base formula.

Induction on \(i\) proves the following statement for all \(m\)-bit blocks,
including invalid ones: \(R_i(u,v)\) is true exactly when there are blocks
\[
 u=w_0,w_1,\ldots,w_\ell=v,
 \qquad \ell\leq2^i,
\]
such that every adjacent pair satisfies either \(\operatorname{Eq}\) or
\(\operatorname{Step}_{M,x}\).  The base is the displayed definition of
\(R_0\).  For the successor case, the existential block \(z\) is a midpoint,
and the universal bit demands both halves: when \(b=0\) the sole recursive
occurrence receives \((u,z)\), and when \(b=1\) it receives \((z,v)\).
Concatenating the two paths proves one direction.  Conversely, on a path of
length at most \(2^{i+1}\), choose \(z\) after
\(\min\{\ell,2^i\}\) edges; both resulting pieces have length at most
\(2^i\), and the selector supplies the required endpoints for each value of
\(b\).

Because \(\operatorname{Step}_{M,x}\) relates only valid blocks, a sequence
beginning at a valid block cannot enter an invalid block: equality preserves
the block, and a step has a valid target.  Hence, for valid \(u,v\), the
induction statement is exactly reachability of the configuration encoded by
\(v\) from that encoded by \(u\) in at most \(2^i\) machine steps; equality
accounts only for unused allowance.

There are at most \(2^m\) encoded configurations.  A deterministic accepting
run reaches the canonical accepting configuration without repeating a
configuration, since a repetition would loop forever.  Therefore, with the
constant blocks \(c_{\mathrm{init}}(x)\) and \(c_{\mathrm{acc}}\),
\[
 \Phi_{M,x}=R_m(c_{\mathrm{init}}(x),c_{\mathrm{acc}})
\]
is true exactly when \(M\) accepts \(x\).  Its size is polynomial because the
base is polynomial and the recurrence adds only \(O(m)\) symbols at each of
\(m\) levels; it is constructed in polynomial time with fresh variable names.
This gives a polynomial-time many-one reduction from every language in
\(\mathsf{PSPACE}\) to \(\mathsf{TQBF}\), completing the two directions used
by Theorem~\ref{thm:p04-tqbf-pspace}.
\end{annexitem}
